\documentclass[10pt,a4paper]{amsart}
\usepackage[margin=19mm]{geometry}
\usepackage{amsmath,amssymb,mathtools}
\usepackage[scr=rsfs,cal=euler]{mathalfa}
\usepackage{xfrac}
\usepackage[unicode,hidelinks,bookmarks=false]{hyperref}
\newcommand{\ie}{i.e.\ }
\newcommand{\cf}{cf.\ }
\newcommand{\resp}{respectively\ }
\newcommand{\Subsection}{\subsection}
\providecommand{\nobreakdash}{\nobreak\hbox{-}}
\setlength{\emergencystretch}{2em}
\multlinegap0pt
\usepackage{amssymb}
\newtheorem{theorem}{Theorem}
\DeclareMathOperator{\Aut}{Aut}
\newcommand{\CV}{\mathcal{V}}
\DeclareMathOperator{\Com}{Com}
\newcommand{\ZZ}{\mathbb{Z}}
\newcommand{\alg}[1]{\mathfrak{#1}}
\newcommand{\be}{\begin{equation}}
\newcommand{\ee}{\end{equation}}
\newcommand{\paraf}{\mathcal{P}}
\title{Monstrous parafermionic defects and other non-invertible symmetries in chiral CFTs}
\author{Roberto Volpato}
\date{Source: arXiv:2609.03043v2 (CC BY 4.0)}
\begin{document}
\maketitle

\noindent\emph{Source and licence.} Monstrous parafermionic defects and other non-invertible symmetries in chiral CFTs. Version arXiv:2609.03043v2 (CC BY 4.0), distributed by arXiv under the Creative Commons Attribution 4.0 International licence, http://creativecommons.org/licenses/by/4.0/, as recorded in the arXiv licence metadata for this exact version and shown on the abstract page https://arxiv.org/abs/2609.03043v2. Full licence notice: \texttt{licenses/arxiv-2609.03043-CC-BY-4.0.txt}.\par\smallskip

\noindent\textit{Editorial scope.} Counted unit: the article's theorem th:paraf (source0.tex lines 331--344). The article's complete original proof of it is delivered with it (source0.tex lines 345--364, one \texttt{proof} environment of 20 lines). No mathematics is written, repaired or re-derived: the statement and the proof are the article's own lines, in the article's own notation, and no retained line is substituted or re-flowed.  Retained uncounted as the source-local context the proof needs: lines 275--282.  Nothing was omitted from any retained span, so the package carries no omission record.  No retained line contains a citation, so the wrapper carries no bibliography and no bibliography entry.  No retained line contains a figure environment, an image-inclusion command or drawing code, so no image is delivered and the package declares no asset dependency.  Editorial formatting only, in addition: an AMS \texttt{amsart} preamble that restates the article's own declarations and macros used by the retained lines, namely the environments \texttt{theorem} on separate counters, together with the standard packages those lines need.  The retained environments print in this document as Theorem 1 in order of appearance, whereas the article prints the same items with the numbering of its own full document.  The complete native source of the article as distributed in the arXiv e-print for this exact version is bundled unchanged as \texttt{sources/209257/source0.tex}.
\begin{theorem}\label{th:simplecurrent}
	Let $V$ be a strongly rational VOA admitting a simple current $J$ of order $N$.  Suppose that $J$ contains a vector $v\neq 0$ such that:
	\begin{enumerate}
		\item the conformal weight of $v$ is $h(v)=1-\frac{1}{N}$;
		\item either $V$ contains no states of weight $1$ ($V_1=0$), or $v$ is in the vacuum representation of the affine Kac-Moody algebra $\alg{h}$ generated by the operators in $V_1$;
	\end{enumerate}
then, there is an embedding of the $\ZZ_N$ parafermion algebra $\paraf(N)=\frac{\alg{su}(2)_N}{\alg{u}(1)_N}$ in $V$.
\end{theorem}

\begin{theorem}\label{th:paraf}
	Let $V$ be a strongly rational holomorphic VOA, and $g\in \Aut(V)$ be a non-anomalous automorphism of finite order $N$. Let $(V^g)_1$ the space of currents (vectors of conformal weight $1$) of  the $g$-fixed subVOA $V^g\subset V$. If the $g$-twisted sector $V_g$ contains a non-zero vector $v\in V_g$ such that:
	\begin{enumerate}
		\item the conformal weight $h(v)$ satisfies
		$$ 0<h(v)<1
		$$ and
		\item either $(V^g)_1=0$, or $v$ is in the vacuum representation of the current algebra $\alg{h}$ generated by the weight $1$ operators in $(V^g)_1$,
	\end{enumerate} then:
	\begin{enumerate}
		\item The conformal weight of $v$ is necessarily $h(v)=1-1/N$;
		\item $V$ is self-orbifold with respect to $g$, i.e. the orbifold VOA $V/\langle g\rangle$ is isomorphic to $V$;
		\item There is an embedding of the $\ZZ_N$ parafermion algebra $\paraf(N)$ in $V^g$. %Furthermore, the group $C(g)/\langle g\rangle$ acts faithfully by automorphisms on the commutant $\Com_{V^g}(\paraf(N))$ of $\paraf(N)$ in $V^g$, where $C(g)=\{h\in\Aut(V)\mid hg=gh\}$ is the centralizer of $g$ in $\Aut(V^g)$.
	\end{enumerate} 
\end{theorem}

\begin{proof}
Because $g$ is not anomalous, the conformal weight of the $g$-twisted sector is in $\frac{1}{N}\ZZ$, and the condition $0<h(v)<1$ implies that $h(v)=1-\frac{r}{N}$ for some integral $r$ with $0<r<N$. Using the decomposition $V_g=\oplus_{b\in \ZZ/N\ZZ} V_{1,b}$ into $g$-eigenspaces, and the fact that $V_{1,b}$ has conformal weights in $\frac{b}{N}+\ZZ$, we conclude that $v$ must be contained in the component $V_{1,-r}$. Furthermore, a state $v'$ in the vacuum representation of the current algebra $\alg{h}$ and with the same conformal weight must be contained in the dual (contragredient) module $V_{-1,r}\subset V_{g^{-1}}$. Let us show that $r$ must be $1$. We  use an analogous construction as in the proof ot theorem \ref{th:simplecurrent}, but now based on the free boson VOA $W_0:=V_{\sqrt{2rN}\ZZ}$ based on the lattice $\sqrt{2rN}\ZZ$. We take the product $W_0\otimes V^g$, and denote by
\be V_{m,a,b}:=W_m\otimes V_{a,b}\ ,\qquad m\in \ZZ/2rN\ZZ,\quad a,b\in \ZZ/N\ZZ\ ,
\ee the modules of this product algebra. Then, we can extend $V_{0,0,0}\cong W_0\otimes V^g$ by the simple current $V_{2r,1,-r}=W_{2r}\otimes V_{1,-r}$ of order $N$, that has a ground state $\CV_{\frac{+2r}{\sqrt{2rN}}}\otimes v$ of conformal weight $1-\frac{r}{N}+\frac{(2r)^2}{4rN}=1$, to get the VOA
\be \tilde V=\oplus_{n\in \ZZ/N\ZZ} V_{2rn, n,-nr}\ .
\ee The modules of $\tilde V$ decompose into modules $V_{m,a,b}$ of the subVOA $V_{0,0,0}$ that are local with respect to the simple current $V_{2r,1,-r}$, i.e. such that the difference of conformal weights $h(V_{m+2r,a+1,b-r})-h(V_{m,a,b})$ is integral. The latter condition translates into
\begin{align} \frac{(m+2r)^2}{4Nr}+\frac{(a+1)(b-r)}{N}- \frac{m^2}{4Nr}-\frac{ab}{N}\in \ZZ \quad &\Leftrightarrow \quad \frac{4r(m+r)}{4Nr}+\frac{b-r-ar}{N}\in \ZZ  \\ &\Leftrightarrow \quad  m-ar+b\equiv 0\mod N\ZZ\ .
\end{align} The modules of $\tilde V$ are therefore given by
\be \tilde V_t=\bigoplus_{n\in \ZZ/N\ZZ} V_{2nr-t,n,t-nj}\qquad t\in \ZZ/2N\ZZ\ .
\ee By the same argument as in the proof of theorem \ref{th:simplecurrent}, the VOA $\tilde V$ contains an affine subalgebra $\alg{su}(2)$ at some level. Therefore, the group of inner automorphisms of $\tilde V$ contains a $SO(3)$ subgroup generated by the zero modes of these $\alg{su}(2)$ currents. In particular, there is an involution $h$ whose adjoint action on the three $\alg{su}(2)$ currents  maps $j^0(z)$ to $-j^0(z)$. Being an inner automorphism, it must map each module $\tilde V_t$ into itself, and it must flip the sign of the $j^0$-charge. This implies that in each module $\tilde V_t$, the set of $j^0$-charges must be symmetric with respect to change of sign. But if we take, for example, the module
\be \tilde V_1=\bigoplus_{n\in \ZZ/N\ZZ} V_{2nr-1,n,1-nr}\ ,
\ee we see that the $j^0$-charges take values in $\frac{1}{\sqrt{2rN}} (-1+2r\ZZ)$. This set is symmetric with respect to $0$ if and only if $r=1$, so this is the only consistent value of $r$, and point (1) of the theorem is proved.\\
Let us now set $r=1$, and prove that the theory is self-orbifold, i.e $V\cong V/\langle g\rangle$. The inner automorphism $h$ maps $j^0$ to $-j^0$, so it must map the $0$-charge component $V_{0,0,0}=W_0\otimes V^g$ of $\tilde V$ to itself, and it must be given by the tensor product $h=f_{W}\otimes f_{V^g}$ of the charge conjugation automorphism $f_W$ in $W_0$ and an involution $f_{V^g}$ of  $\Com(\alg{u}(1),\tilde V)\cong V^g$. As mentioned above, because it is inner it must map each $\tilde V$-module 
\be \tilde V_t=\bigoplus_{n\in \ZZ/N\ZZ} V_{2n-t,n,t-n}=\bigoplus_{\substack{m\in \ZZ/2N\ZZ\\m\equiv t\mod 2}} V_{m,\frac{m+t}{2},\frac{t-m}{2}}=\bigoplus_{\substack{m\in \ZZ/2N\ZZ\\m\equiv t\mod 2}} W_m\otimes V_{\frac{m+t}{2},\frac{t-m}{2}}
\ee  into itself, and map each component $ V_{m,\frac{m+t}{2},\frac{t-m}{2}}$ to the component $V_{-m,\frac{-m+t}{2},\frac{t+m}{2}}$ with opposite $j^0$ charge. As a consequence, the involution $f_{V^g}$ must map each $V^g$-module $V_{a,b}$ to the module $V_{b,a}$, for all $a,b\in \ZZ/N\ZZ$. Because $V$ and $V/\langle g\rangle$ are given by
\be V=\bigoplus_{n\in \ZZ/N\ZZ} V_{0,n}\ ,\qquad V/\langle g\rangle=\bigoplus_{n\in \ZZ/N\ZZ} V_{n,0}
\ee it follows that the involution $f_{V^g}$ extends to a VOA isomorphism $f: V\to V/\langle g\rangle$, thus proving point (2).
\\
Finally, to prove point (3) we just notice that, because $r=1$, then the hypotheses of theorem \ref{th:simplecurrent} are all satisfied, so that there is an embedding of the parafermion algebra $\paraf(N)$ in $V^g$.
\end{proof}

\end{document}
